% from ega4-13-fr.tex
\oldpage[IV-3]{189}
caténaire (5.6.4)~; par suite, pour tout $y_0\in f_0(X_0)$, toutes les
composantes irréductibles de $f_0^{-1}(y)$ ont une dimension $\geq e$, car on
sait que $\dim(f_0^{-1}(\eta_0))=e$ par transitivité des fibres (4.1.4). Pour
tout $y\in f(X)$, on a alors $y_0=g(y)\in f_0(X_0)$ et par transitivité des
fibres et (4.2.7), on achève la démonstration.
\end{proof}

\begin{theorem}[13.1.3]
\label{IV.13.1.3-fr}
\textnormal{(Chevalley).} --- Soit $f:X\to Y$ un morphisme localement de type
fini. Pour tout entier $n$, l'ensemble $F_n(X)$ des $x\in X$ tels que
$\dim_x(f^{-1}(f(x)))\geq n$ est fermé~; autrement dit, la fonction
$x\mapsto\dim_x(f^{-1}(f(x)))$ est semi-continue supérieurement dans $X$.
\end{theorem}

\begin{proof}
I) Supposons d'abord que $f$ soit localement de présentation finie.

La question est évidemment locale sur $X$ et sur $Y$, et on peut donc supposer
$Y=\operatorname{Spec}(A)$, $X=\operatorname{Spec}(B)$ affines, $B$ étant une
$A$-algèbre de présentation finie. On sait alors (8.9.1) qu'il y a un
sous-anneau noethérien $A_0$ de $A$, et une $A_0$-algèbre de type fini $B_0$
tels que si l'on pose $Y_0=\operatorname{Spec}(A_0)$,
$X_0=\operatorname{Spec}(B_0)$, on ait

\[
  X=X_0\times_{Y_0}Y,\qquad f=f_0\times 1_Y,
\]

où $f_0:X_0\to Y_0$ correspond à l'homomorphisme $A_0\to B_0$. Soient
$g:Y\to Y_0$ le morphisme correspondant à l'injection canonique $A_0\to A$,
$y$ un point de $Y$, $y_0=g(y)$~; on sait que
$f^{-1}(y)=f_0^{-1}(y_0)\otimes_{k(y_0)}k(y)$, et il résulte de (4.2.7) que
si $p$ est la projection canonique $f^{-1}(y)\to f_0^{-1}(y_0)$, les
composantes irréductibles de $f^{-1}(y)$ sont les composantes irréductibles des
ensembles $p^{-1}(Z_0)$, où $Z_0$ parcourt l'ensemble des composantes
irréductibles de $f_0^{-1}(y_0)$, et chacune des composantes irréductibles de
$p^{-1}(Z_0)$ domine $Z_0$ et a une dimension égale à $\dim(Z_0)$. Tenant
compte de ($\mathbf 0$, 14.1.5), on voit donc que si $g':X\to X_0$ est la
projection canonique, $x$ un point de $X$ et $x_0=g'(x)$, on a
$\dim_x(f^{-1}(f(x)))=\dim_{x_0}(f_0^{-1}(f_0(x_0)))$, d'où
$F_n(X)=g'^{-1}(F_n(X_0))$, et on est par suite ramené à démontrer le théorème
lorsque $Y$ est \emph{noethérien}, ce que nous supposerons désormais.

On peut évidemment supposer $X$ et $Y$ réduits (1.5.4). En considérant
l'ensemble des parties fermées $Y'$ de $Y$ telles que le théorème soit vrai
pour le sous-préschéma fermé de $Y$ ayant $Y'$ pour espace sous-jacent et pour
$X'=f^{-1}(Y')$, on peut raisonner par récurrence noethérienne
($\mathbf 0_{\mathrm I}$, 2.2.2), et supposer que le théorème est vrai pour
toute partie fermée $Y'\neq Y$ de $Y$. Si $X_i$ ($1\leq i\leq m$) sont les
sous-préschémas fermés réduits de $X$ ayant pour espaces sous-jacents les
composantes irréductibles de $X$, on a $F_n(X)=\bigcup_iF_n(X_i)$ en vertu de
($\mathbf 0$, 14.1.5), et l'on peut donc se borner à démontrer le théorème pour
chacun des $X_i$, autrement dit, on peut supposer $X$ irréductible. Si $Z$ est
le sous-préschéma fermé de $Y$ ayant $\overline{f(X)}$ pour espace sous-jacent,
$f$ se factorise en $X\xrightarrow{g}Z\xrightarrow{j}Y$, où $j$ est
l'injection canonique (\textbf{I}, 5.2.2), et $g$ est de type fini (1.5.4),
donc il suffit de démontrer le théorème pour $Z$ et $g$~; en vertu de
l'hypothèse de récurrence, on est donc ramené à considérer seulement le cas où
$Z=Y$, autrement dit où $Y$ est irréductible et $f$ dominant. Soit alors
$\eta$ le point générique de $Y$ et posons $e=\dim(f^{-1}(\eta))$~; il résulte
de (13.1.1) que pour $n\leq e$ on a $F_n(X)=X$, et par suite on peut se borner
au cas où $n>e$. Mais alors (9.5.6), il y a un voisinage ouvert $U$ de $\eta$
dans $Y$ tel que $F_n(X)\subset f^{-1}(Y-U)$~; comme $Y-U\neq Y$, l'hypothèse
de récurrence entraîne que $F_n(X)$ est fermé.

II) Passons maintenant au cas général, en supposant toujours
$S=\operatorname{Spec}(A)$ et $X=\operatorname{Spec}(B)$ affines, $B$ étant
une $A$-algèbre de type fini, donc de la forme

